\section{Threat Model}
\label{sec:threat-model}

The attacker is computationally bounded by the IND-CCA2 security of ML-KEM-1024 \citep{fips203,hofheinz2017fo} and the confidentiality and integrity of the symmetric AEAD (AES-256-GCM under its standard assumptions, augmented with the UtC commitment tag \citep{bellare2022commitment}). They see every byte of every envelope, may control delivery (replay, reorder, drop), may have a decryption oracle for \emph{other} ciphertexts under the same recipient key (IND-CCA2 required), and may flip any header byte or substitute associated data (the AEAD's AD path authenticates both).

The attacker \textbf{cannot} compromise the recipient's long-term ML-KEM-1024 private key (and, for hybrid suite \code{0x06}, the X25519 long-term scalar): that is the plugin protocol's responsibility (\cref{sec:plugin-protocol}). The attacker may attempt \emph{multi-recipient/AEAD-key confusion} \citep{grubbs2017aead,albertini2020aead,len2021aead}; the 32-byte UtC commitment tag closes this class on every suite, including hybrid.

The attacker may attempt \emph{side-channel attacks} against ML-KEM: KyberSlash 1/2 timing leaks \citep{kyberslash}, PQShield power-side-channel disclosures \citep{pqshield2024}, and SPA-CCA \citep{guo2024spa}. They may attempt timing or power side-channels against the X25519 half of suite \code{0x06}. \pqcrypt{} mitigates software-timing leakage by pinning verified, constant-time third-party ML-KEM-1024 and X25519 implementations per language (\cref{sec:mlkem,sec:hybrid}; the only crypto primitive the monorepo \emph{vendors} is Keccak/SHAKE-256, under \code{common/}) and documents an explicit ``constant-time $\neq$ constant-power'' disclaimer (\cref{sec:side-channels}). Physical-access attacks (power, EM, frequency-scaling \citep{hertzbleed}, prefetcher \citep{gofetch}) remain out of v1.x scope.

The attacker may attempt a \emph{suite-confusion} attack: present an envelope with a suite byte that does not match the recipient's key, or re-tag a hybrid envelope as non-hybrid (or vice versa). \pqcrypt{} rejects this structurally at the API boundary (\cref{sec:failure-modes}): \code{encrypt} rejects a recipient key whose suite does not match the requested suite; \code{decrypt} rejects an envelope whose header suite does not match the private key's suite. Both are \texttt{BadKey}-class structural errors, distinguishable safely because they do not depend on secret-correlated bytes (enforced in all eight reference implementations as of 2026-06-10).

\subsection*{Residual risk and the v1.2 hybrid escape valve}

\pqcrypt{}'s \textbf{core promise} (suites \code{0x01}/\code{0x02}) is post-quantum-only by deliberate scope. If ML-KEM-1024 is later structurally broken (a sub-exponential Module-LWE attack), every \code{0x01}/\code{0x02} envelope loses confidentiality with no classical fallback. This 30--50 year risk is accepted in exchange for an unambiguous PQ-only promise; the mitigation path is forward-only via re-encryption under a future suite (\cref{sec:appendices}).

For callers who cannot accept that residual risk, v1.2 introduced the \emph{opt-in} hybrid suite \code{0x06} (\cref{sec:hybrid}). Suite \code{0x06} is IND-CCA2-secure as long as \emph{either} ML-KEM-1024 \emph{or} X25519 is IND-CCA2-secure; an attacker must break both to recover the session key. This is a separate suite allocation, not a modification of the core promise --- the wire is unambiguous about which guarantee a given envelope makes.
